% ch3_c (书页 119-130 / p134-p145)
%--B-TAIL--
%JOIN%
以及
\begin{equation*}
\Gamma^1_{\mu\lambda} T^{\mu\lambda} = \Gamma^1_{00} T^{00} + \Gamma^1_{11} T^{11} + \Gamma^1_{22} T^{22} + \Gamma^1_{33} T^{33} = 0. \tag{3.103}
\end{equation*}

对 $\nu=2$ 和 $\nu=3$ 也会得到等价的结果。于是，能量-动量守恒方程的空间分量简单地化为

\begin{equation*}
\partial_i p = 0. \tag{3.104}
\end{equation*}

把这些结果与我们在 Minkowski 时空中会得到的结果作比较是很有启发性的，后者只需令 $a=1$、$\dot a =0$ 即可得到。压强梯度方程 (3.104) 不受影响，所以曲率对空间分量没有影响：对于一个在共动观者看来静止的流体，压强必须在整个空间中为常数。而时间分量则不同：宇宙的膨胀给 (3.100) 的右端引入了一个非零的项。为了理解这个新特点的后果，让我们考虑如下形式的物态方程

\begin{equation*}
p = w\rho, \tag{3.105}
\end{equation*}

其中 $w$ 是某个常数。那么 (3.100) 变为

\begin{equation*}
\frac{\dot\rho}{\rho} = -3(1+w)\frac{\dot a}{a}, \tag{3.106}
\end{equation*}

求解可得

\begin{equation*}
\rho \propto a^{-3(1+w)}. \tag{3.107}
\end{equation*}

在第 1 章中我们提到过三种物态方程形如 (3.105) 的理想流体：尘埃，$w=0$；辐射，$w=1/3$；以及真空，$w=-1$。一组非相对论性、无相互作用的粒子表现得像尘埃；一组光子或其他无质量粒子表现得像辐射；而遍布时空的非零常数能量密度则表现得像真空。由 (3.107) 我们看到，物态方程决定了能量密度随宇宙膨胀的演化方式：

\begin{align*}
\text{物质}\quad & p = 0 & \rho &\propto a^{-3}\\
\text{辐射}\quad & p = \tfrac{1}{3}\rho & \rho &\propto a^{-4}\\
\text{真空}\quad & p = -\rho & \rho &= \text{常数}. \tag{3.108}
\end{align*}

我们将在第 8 章更深入地探讨这些行为；现在只须注意它们是合理的。对尘埃而言，能量密度来自每个粒子的静止质量；如果所有粒子都具有质量 $m$，能量密度就是 $\rho = nm$，其中 $n$ 是数密度（number density）。由于数密度按 $a^{-3}$ 下降（共动区域的物理体积增大，而粒子总数保持不变），同时质量不变，我们预期能量密度满足 $\rho \propto a^{-3}$。对辐射而言，每个粒子（比如光子）的能量随宇宙膨胀按 $a^{-1}$ 红移掉；由于数密度仍满足 $n \propto a^{-3}$，我们预期 $\rho \propto a^{-4}$。最后，真空能量密度是任一物理体积中固有的、不随时间改变的能量；它在宇宙膨胀时完全不红移，所以我们得到 $\rho = \text{常数}$。

这个例子生动地展现了平直时空与弯曲时空之间的差异。例如，考虑一下我们可能禁不住想称之为“能量”的东西，即能量密度对空间的积分：$E = \int \rho a^3\, d^3x$，其中边界取在固定的共动坐标处，因此该区域随宇宙一起膨胀，而因子 $a^3$ 来自空间度规 $a^2\delta_{ij}$ 的行列式的平方根。这个量一般而言显然不守恒。对尘埃，由于 $\rho \propto a^{-3}$，$E$ 在宇宙膨胀时保持不变；但对辐射它不断减小，而对真空能量它不断增大。这令人沮丧，因为能量守恒是物理学中最受珍视的原理之一。发生了什么？一种思考方式是从诺特定理（Noether's theorem）出发，该定理指出每一种对称性都蕴含一个守恒量。能量正是由时间平移不变性导出的守恒量。显然，在膨胀宇宙中，能量-动量张量定义在一个随时间变化的背景上；因此没有理由相信能量应当守恒。（用 3.8 节将要引入的语言来说，就是“不存在类时 Killing 矢量”。）尽管如此，我们仍继续把 $\nabla_\mu T^{\mu\nu} = 0$ 称为能量-动量守恒方程。它传达的观念是：能量-动量张量服从一条确定的定律，即使不存在与守恒能量相对应的积分。从平直时空到弯曲时空的转变引入了 (3.96) 中额外的 Christoffel 符号项；粗略地说，这些项的作用是允许能量在物质场（即 $T^{\mu\nu}$）与引力场之间转移。然而这个说法并不十分正式，你也不要把它推得太远——事实证明，很难给引力场联系上一个局域的能量密度，尽管在某些情形下这是可能的。

当然，还有一种时间平移不变性的概念，它指的不是背景时空，而是理论本身（也就是说，指的是定义理论的方程，而不是理论的某个具体解）。我们尚未建立起广义相对论的动力学方程，但它们将被证明在时间平移下不变，而且在任何其他类型的坐标变换下也不变——它们本来就必须如此。这种广义坐标不变性给理论的允许组态加上了一组约束，而且通常需要更精细的分析来处理。

最后，你应当接受这样一点：平直时空与弯曲时空之间存在着深刻的差异，我们来自平直时空物理学的一些心仪观念，在这个更一般的框架中将被大幅修正。这并不是广义相对论有什么缺陷的信号，而是抛弃我们习以为常的刚性时空几何之后的一个自然结果。

\section{Riemann 曲率张量}

在建立了协变导数和平行移动这套机制之后，我们终于准备好正面对待曲率了。曲率由 Riemann 张量来量化，而 Riemann 张量由联络导出。这种曲率量度背后的想法是：我们知道联络的“平直性”是什么意思——与欧几里得型或 Minkowski 型度规相联系的常规（且通常是隐含的）Christoffel 联络具有若干性质，它们可以被看作平直性的不同表现。这包括：沿闭合回路一周的平行移动不改变矢量；张量的协变导数相互对易；初始平行的测地线保持平行。正如我们将看到的，当我们研究这些性质在更一般的情形下如何被改变时，Riemann 张量就出现了。

我们前面已经论证过（以二维球面为例），在弯曲空间中沿闭合回路一周的平行移动会导致矢量发生一个变换。所得的变换依赖于回路所包围的总曲率；更有用的应当是每一点曲率的局域描述，而这正是 Riemann 张量所要提供的。因此，引入 Riemann 张量的一种常规做法是考虑绕无穷小回路的平行移动。这里我们不打算那样做，而是采取一条更直接的路线。不过，即使不细究细节，也能够看出答案应当取什么形式。由于时空在足够小的区域内看起来是平直的，我们的回路将由两个（无穷小的）矢量 $A^\mu$ 和 $B^\nu$ 确定。我们想象把矢量 $V^\mu$ 平行移动：先沿 $A^\mu$ 的方向移动，再沿 $B^\nu$，然后沿 $A^\mu$ 和 $B^\nu$ 逆向退回出发点，如图 3.5 所示。

%FIG{3.5}: {由两个矢量 $A^\mu$ 与 $B^\nu$ 定义的无穷小回路}

我们知道平行移动的作用与坐标无关，所以应当存在某个张量，它告诉我们当矢量回到出发点时发生了怎样的变化；它是作用在矢量上的一个线性变换，因而含有一个上指标和一个下指标。但它还依赖于定义回路的两个矢量 $A$ 和 $B$；因此还应当有两个附加的下指标，用来与 $A^\mu$ 和 $B^\nu$ 缩并。此外，这个张量在这两个指标上应当是反对称的，因为交换这两个矢量相当于沿相反方向绕行回路，应当给出原答案的逆。这与“若 $A$ 与 $B$ 是同一个矢量则变换应为零”的事实一致。于是我们预期，这个矢量绕回路平行移动一周所经历的改变 $\delta V^\rho$ 应当具有如下形式

\begin{equation*}
\delta V^\rho = R^\rho{}_{\sigma\mu\nu} V^\sigma A^\mu B^\nu, \tag{3.109}
\end{equation*}

其中 $R^\rho{}_{\sigma\mu\nu}$ 是一个 $(1,3)$ 型张量，称为 \textbf{Riemann 张量}（Riemann tensor，或简称曲率张量）。它在后两个指标上反对称：

\begin{equation*}
R^\rho{}_{\sigma\mu\nu} = -R^\rho{}_{\sigma\nu\mu}. \tag{3.110}
\end{equation*}

当然，如果把 (3.109) 取作 Riemann 张量的定义，就需要为指标的排列顺序选定一个约定。关于这个约定应当是什么，完全没有共识，所以要当心。

基于我们对平行移动的了解，本可以非常仔细地完成必要的操作，看看矢量在这一操作下会发生什么，其结果将是一个用联络系数表示的曲率张量公式。然而，快得多的办法是考虑一个相关的操作，即两个协变导数的对易子。它与绕回路平行移动之间的关系应当一目了然：张量沿某一方向的协变导数量度的是，该张量相对于它被平行移动时应有的样子变化了多少，因为张量沿它被平行移动的那个方向的协变导数为零。于是，两个协变导数的对易子量度的就是“先把张量沿一个方向平行移动、再沿另一个方向”与“次序相反”这两种做法之间的差别，如图 3.6 所示。

%FIG{3.6}: {两个协变导数的对易子}

实际计算非常直接。考虑矢量场 $V^\rho$，我们取

\begin{align*}
[\nabla_\mu, \nabla_\nu] V^\rho &= \nabla_\mu\nabla_\nu V^\rho - \nabla_\nu\nabla_\mu V^\rho\\
&= \partial_\mu(\nabla_\nu V^\rho) - \Gamma^\lambda_{\mu\nu}\nabla_\lambda V^\rho + \Gamma^\rho_{\mu\sigma}\nabla_\nu V^\sigma - (\mu \leftrightarrow \nu)\\
&= \partial_\mu\partial_\nu V^\rho + (\partial_\mu\Gamma^\rho_{\nu\sigma})V^\sigma + \Gamma^\rho_{\nu\sigma}\partial_\mu V^\sigma - \Gamma^\lambda_{\mu\nu}\partial_\lambda V^\rho - \Gamma^\lambda_{\mu\nu}\Gamma^\rho_{\lambda\sigma}V^\sigma\\
&\quad + \Gamma^\rho_{\mu\sigma}\partial_\nu V^\sigma + \Gamma^\rho_{\mu\sigma}\Gamma^\sigma_{\nu\lambda}V^\lambda - (\mu \leftrightarrow \nu)\\
&= (\partial_\mu\Gamma^\rho_{\nu\sigma} - \partial_\nu\Gamma^\rho_{\mu\sigma} + \Gamma^\rho_{\mu\lambda}\Gamma^\lambda_{\nu\sigma} - \Gamma^\rho_{\nu\lambda}\Gamma^\lambda_{\mu\sigma})V^\sigma - 2\Gamma^\lambda_{[\mu\nu]}\nabla_\lambda V^\rho. \tag{3.111}
\end{align*}

在最后一步中，我们对一些哑指标做了改名，并删去了反对称化后相消的一些项。我们认出，最后一项中被反对称化的联络系数恰好是挠率张量的一半，而左端显然是张量；因此括号中的表达式本身必定是张量。我们写成

\begin{equation*}
[\nabla_\mu, \nabla_\nu] V^\rho = R^\rho{}_{\sigma\mu\nu}V^\sigma - T^\lambda{}_{\mu\nu}\nabla_\lambda V^\rho, \tag{3.112}
\end{equation*}

其中 Riemann 张量被确认为

\begin{equation*}
R^\rho{}_{\sigma\mu\nu} = \partial_\mu\Gamma^\rho_{\nu\sigma} - \partial_\nu\Gamma^\rho_{\mu\sigma} + \Gamma^\rho_{\mu\lambda}\Gamma^\lambda_{\nu\sigma} - \Gamma^\rho_{\nu\lambda}\Gamma^\lambda_{\mu\sigma}. \tag{3.113}
\end{equation*}

关于这个表达式的推导，有几点值得注意：

\begin{itemize}
\item 当然，我们并没有证明 (3.113) 与 (3.109) 中出现的确实是同一个张量，但事实上确实如此。习题中会请你证明这一点。
\item 或许令人惊讶的是，对易子 $[\nabla_\mu, \nabla_\nu]$ 虽然看似一个微分算符，但它作用在矢量场上的效果（至少在没有挠率时）是一个简单的乘法变换。Riemann 张量量度的是协变导数对易子中正比于矢量场的那一部分，而挠率张量量度的是正比于矢量场的协变导数的那一部分；二阶导数完全没有出现。
\item 注意 (3.113) 是由非张量性的元素构造出来的；你可以验证，各个变换定律恰好都行得通，使这个特定的组合成为一个合法的张量。
\item $R^\rho{}_{\sigma\mu\nu}$ 在其最后两个指标上的反对称性可以直接从这个公式及其推导看出。
\item 我们完全由联络构造了曲率张量（完全没有提及度规）。我们足够小心，因此上面的表达式对任何联络都成立，无论它是否度规相容或无挠。
\item 用我们如今已惯用的方法，可以计算出 $[\nabla_\rho, \nabla_\sigma]$ 作用在任意阶张量上的结果。答案是
\begin{align*}
[\nabla_\rho, \nabla_\sigma] X^{\mu_1\cdots\mu_k}{}_{\nu_1\cdots\nu_\ell}
&= -T^\lambda{}_{\rho\sigma}\nabla_\lambda X^{\mu_1\cdots\mu_k}{}_{\nu_1\cdots\nu_\ell}\\
&\quad + R^{\mu_1}{}_{\lambda\rho\sigma}X^{\lambda\mu_2\cdots\mu_k}{}_{\nu_1\cdots\nu_\ell} + R^{\mu_2}{}_{\lambda\rho\sigma}X^{\mu_1\lambda\cdots\mu_k}{}_{\nu_1\cdots\nu_\ell} + \cdots\\
&\quad - R^\lambda{}_{\nu_1\rho\sigma}X^{\mu_1\cdots\mu_k}{}_{\lambda\nu_2\cdots\nu_\ell} - R^\lambda{}_{\nu_2\rho\sigma}X^{\mu_1\cdots\mu_k}{}_{\nu_1\lambda\cdots\nu_\ell} - \cdots. \tag{3.114}
\end{align*}
\end{itemize}

挠率张量与 Riemann 张量都可以看作多重线性映射，用矢量场对易子可以把它们表达得十分优雅。把挠率看作从两个矢量场到第三个矢量场的映射，我们有

\begin{equation*}
T(X, Y) = \nabla_X Y - \nabla_Y X - [X, Y], \tag{3.115}
\end{equation*}

而把 Riemann 张量看作从三个矢量场到第四个矢量场的映射，我们有（这套记号看上去有点怪，但它是标准的）

\begin{equation*}
R(X, Y)Z = \nabla_X\nabla_Y Z - \nabla_Y\nabla_X Z - \nabla_{[X, Y]}Z. \tag{3.116}
\end{equation*}

在这些表达式中，记号 $\nabla_X$ 指沿矢量场 $X$ 的协变导数；用分量写即 $\nabla_X = X^\mu\nabla_\mu$。举例来说，(3.116) 等价于

\begin{align*}
R^\rho{}_{\sigma\mu\nu} X^\mu Y^\nu Z^\sigma &= X^\lambda\nabla_\lambda(Y^\eta\nabla_\eta Z^\rho) - Y^\lambda\nabla_\lambda(X^\eta\nabla_\eta Z^\rho)\\
&\quad - (X^\lambda\partial_\lambda Y^\eta - Y^\lambda\partial_\lambda X^\eta)\nabla_\eta Z^\rho, \tag{3.117}
\end{align*}

你可以验证它与 (3.113) 等价。注意 (3.116) 中的两个矢量 $X$ 和 $Y$ 对应于 Riemann 张量分量形式中的最后两个指标。(3.116) 中含对易子 $[X, Y]$ 的最后一项，在 $X$ 和 $Y$ 取坐标基矢量场时为零（因为 $[\partial_\mu, \partial_\nu] = 0$），这正是我们当初计算两个协变导数的对易子时这一项没有出现的原因。我们不会大量使用这套记号，但你可能会在文献中遇到它，所以应当能够解读它。

在把曲率张量定义为刻画联络的量之后，现在应当承认：在广义相对论中我们最关心的是 Christoffel 联络。此时联络由度规导出，相应的曲率可以视为度规本身的曲率。这一对应使我们终于能够理解那个非正式的观念：度规看起来是欧几里得型或 Minkowski 型的空间是平直的。事实上，这在两个方向上都成立：

\begin{itemize}
\item 如果存在一个坐标系，其中度规的分量是常数，那么 Riemann 张量将为零。
\item 如果 Riemann 张量为零，我们总可以构造一个坐标系，使度规分量是常数。
\end{itemize}

严格说来，这些论断应当限制在流形的单连通（simply-connected）区域上（区域内的所有回路都能在不离开该区域的前提下光滑地形变为一点）；下面我们将默认这个条件。

第一条很容易证明。如果我们所在的坐标系使 $\partial_\sigma g_{\mu\nu} = 0$ 处处成立，而不仅仅在一点成立，那么 $\Gamma^\lambda_{\mu\nu} = 0$ 且 $\partial_\sigma\Gamma^\lambda_{\mu\nu} = 0$；于是由 (3.113) 得 $R^\rho{}_{\sigma\mu\nu} = 0$。但这是张量方程，只要它在某个坐标系中成立，它在任何坐标系中都必然成立。因此，“Riemann 张量为零”是“有可能找到使 $g_{\mu\nu}$ 的分量处处为常数的坐标”的必要条件。

第二个论断——$R^\rho{}_{\sigma\mu\nu} = 0$ 处处成立蕴含我们能找到一个使度规分量处处为常数的坐标系——证明起来更难（但也难不到哪里去）。作为热身，考虑在某点 $p$ 定义的一个 1-形式 $\omega = \omega_\mu dx^\mu$。对每条包含 $p$ 的路径 $x^\mu(\lambda)$，只要要求 $\omega_\mu$ 被平行移动，我们就能沿该路径构造一个唯一的 1-形式场：

\begin{equation*}
\frac{dx^\mu}{d\lambda}\nabla_\mu\omega_\nu = 0. \tag{3.118}
\end{equation*}

一般而言，如果我们对从 $p$ 出发、经过另一点 $q$ 的不同路径执行这一步骤，$\omega_\mu$ 在 $q$ 处的值将依赖于路径。然而，若 Riemann 张量处处为零，平行移动便与路径无关，我们就能在整个流形上定义一个唯一的 1-形式场。因此 (3.118) 必须对任意的 $dx^\mu/d\lambda$ 都成立；这只有在 $\omega_\mu$ 为协变常数时才可能：

\begin{equation*}
\nabla_\mu\omega_\nu = 0. \tag{3.119}
\end{equation*}

在任意流形上这个方程一般无解；它在这里有解，只是因为我们假设了曲率为零。我们可以取 (3.119) 的反对称部分，而由 (3.36) 可知这正是外微分：

\begin{equation*}
\nabla_{[\mu}\omega_{\nu]} = \partial_{[\mu}\omega_{\nu]} = 0, \tag{3.120}
\end{equation*}

或者，用无指标记号写成

\begin{equation*}
d\omega = 0. \tag{3.121}
\end{equation*}

换言之，$\omega$ 是闭的。它同时也是恰当的（存在一个标量函数 $\alpha$ 使得 $\omega = d\alpha$），因为我们已经限制了所工作区域的拓扑。用分量表示即

\begin{equation*}
\omega_\mu = \partial_\mu\alpha. \tag{3.122}
\end{equation*}

1-形式 $\omega$ 并没有任何特殊之处，所以我们可以在 $n$ 维流形上用一组 1-形式 $\hat\theta^{(a)}$（其中 $a \in \{1 \ldots n\}$）重复这一步骤。我们可以选取这组 1-形式，使它们构成对偶空间 $T^*_p$ 的一个归一化基，使得度规相对于这组基的分量就是规范形式的分量；换句话说，

\begin{equation*}
ds^2(p) = \eta_{ab}\, \hat\theta^{(a)}\otimes\hat\theta^{(b)}. \tag{3.123}
\end{equation*}

这里我们在广义的意义上使用 $\eta_{ab}$：一个对角元各自为 $+1$ 或 $-1$、其余元素为零的矩阵。$+1$ 与 $-1$ 的具体排列取决于度规的规范形式，但与当前的论证无关。现在让我们把整组基 1-形式平行移动到流形各处；Riemann 张量的消失保证了结果与所走的路径无关。由于相对于度规相容的联络，度规总是被自动地平行移动，度规分量将保持不变，

\begin{equation*}
ds^2(\text{任何地方}) = \eta_{ab}\, \hat\theta^{(a)}\otimes\hat\theta^{(b)}. \tag{3.124}
\end{equation*}

于是我们指定了一组 1-形式场，它们处处定义了一个使度规分量为常数的基。这完全没有什么了不起；无论曲率是多少，在任何流形上都能做到。我们想要表明的是，这是一个\emph{坐标}基（只有当曲率消失时这才可能）。然而，用导出 (3.122) 的同样论证可以知道，所有 $\hat\theta^{(a)}$ 都是恰当形式，因此存在一组函数 $y^a$，使这些 1-形式场是它们的梯度，

\begin{equation*}
\hat\theta^{(a)} = dy^a. \tag{3.125}
\end{equation*}

这 $n$ 个函数正是我们寻求的坐标；在整个流形上，度规取如下形式

\begin{equation*}
ds^2 = \eta_{ab}\, dy^a dy^b. \tag{3.126}
\end{equation*}

到了这里，如果你愿意，尽可以把指标从 $a$、$b$ 换回 $\mu$、$\nu$。

这样我们就验证了：Riemann 张量为我们提供了一个判据，可以回答“某个面目可憎的度规是否其实是别扭坐标系中平直空间的度规”这个问题。如果我们算出这种度规的 Riemann 张量并发现它为零，就知道度规是平直的；若它不为零，就存在曲率。

\section{Riemann 张量的性质}

Riemann 张量有四个指标，在 $n$ 维空间中天真地数有 $n^4$ 个独立分量。实际上，反对称性质 (3.110) 意味着这最后两个指标只能取 $n(n-1)/2$ 个独立值，于是剩下 $n^3(n-1)/2$ 个独立分量。然而当我们考虑 Christoffel 联络时，另外一些对称性会把独立分量的数目进一步削减。现在就来考察这些对称性。

推导这些附加对称性的最简单办法，是考察全下标形式的 Riemann 张量

\begin{equation*}
R_{\rho\sigma\mu\nu} = g_{\rho\lambda}R^\lambda{}_{\sigma\mu\nu}. \tag{3.127}
\end{equation*}

让我们进一步考虑这个张量在点 $p$ 处建立的局部惯性坐标 $x^{\hat\mu}$ 中的分量。此时 Christoffel 符号本身将为零，但它们的导数不为零。于是我们有

\begin{align*}
R_{\hat\rho\hat\sigma\hat\mu\hat\nu}(p) &= g_{\hat\rho\hat\lambda}\big(\partial_{\hat\mu}\Gamma^{\hat\lambda}{}_{\hat\nu\hat\sigma} - \partial_{\hat\nu}\Gamma^{\hat\lambda}{}_{\hat\mu\hat\sigma}\big)\\
&= \tfrac{1}{2} g_{\hat\rho\hat\lambda} g^{\hat\lambda\hat\tau}\big(\partial_{\hat\mu}\partial_{\hat\nu}g_{\hat\sigma\hat\tau} + \partial_{\hat\mu}\partial_{\hat\sigma}g_{\hat\tau\hat\nu} - \partial_{\hat\mu}\partial_{\hat\tau}g_{\hat\nu\hat\sigma} - \partial_{\hat\nu}\partial_{\hat\mu}g_{\hat\sigma\hat\tau}\\
&\quad - \partial_{\hat\nu}\partial_{\hat\sigma}g_{\hat\tau\hat\mu} + \partial_{\hat\nu}\partial_{\hat\tau}g_{\hat\mu\hat\sigma}\big)\\
&= \tfrac{1}{2}\big(\partial_{\hat\mu}\partial_{\hat\sigma}g_{\hat\rho\hat\nu} - \partial_{\hat\mu}\partial_{\hat\rho}g_{\hat\nu\hat\sigma} - \partial_{\hat\nu}\partial_{\hat\sigma}g_{\hat\rho\hat\mu} + \partial_{\hat\nu}\partial_{\hat\rho}g_{\hat\mu\hat\sigma}\big). \tag{3.128}
\end{align*}

第一行我们用了 $\Gamma^{\hat\tau}{}_{\hat\mu\hat\nu}(p) = 0$，第二行用了 Riemann 法坐标（Riemann normal coordinates）中的 $\partial_{\hat\mu}g^{\hat\lambda\hat\tau} = 0$，第三行则用了偏导数可对易这一事实。从这个表达式出发，我们可以立刻注意到 $R_{\rho\sigma\mu\nu}$ 的三个性质：它在头两个指标上反对称，

\begin{equation*}
R_{\rho\sigma\mu\nu} = -R_{\sigma\rho\mu\nu}, \tag{3.129}
\end{equation*}

在后两个指标上反对称［这一点我们已经从 (3.110) 知道了］，

\begin{equation*}
R_{\rho\sigma\mu\nu} = -R_{\rho\sigma\nu\mu}, \tag{3.130}
\end{equation*}

并且把头一对指标与后一对指标互换时它不变：

\begin{equation*}
R_{\rho\sigma\mu\nu} = R_{\mu\nu\rho\sigma}. \tag{3.131}
\end{equation*}

稍微再多做一点工作（细节留给你想象），可以看出最后三个指标的循环置换（cyclic permutations）之和为零：

\begin{equation*}
R_{\rho\sigma\mu\nu} + R_{\rho\mu\nu\sigma} + R_{\rho\nu\sigma\mu} = 0. \tag{3.132}
\end{equation*}

给定 (3.130)，容易看出这最后一条性质等价于最后三个指标的反对称部分为零：

\begin{equation*}
R_{\rho[\sigma\mu\nu]} = 0. \tag{3.133}
\end{equation*}

所有这些性质都是在特殊坐标系中导出的，但它们都是张量方程，因此在任何坐标中都成立（所以我们也就懒得给指标加帽子了）。它们并非全都独立；只要花些功夫，你可以证明 (3.129)、(3.130) 和 (3.133) 合起来蕴含 (3.131)。这些方程之间的逻辑依存关系，通常不如“它们为真”这一事实来得重要。

知道了 Riemann 张量不同分量之间的这些关系，还剩下多少个独立的量？让我们从如下事实出发：$R_{\rho\sigma\mu\nu}$ 在头两个指标上反对称，在后两个指标上反对称，并且交换这两对指标时对称。这意味着我们可以把它看作对称矩阵 $R_{[\rho\sigma][\mu\nu]}$，其中指标对 $\rho\sigma$ 和 $\mu\nu$ 被当作单个指标。一个 $m\times m$ 对称矩阵有 $m(m+1)/2$ 个独立分量，而一个 $n\times n$ 反对称矩阵有 $n(n-1)/2$ 个独立分量。于是我们得到

\begin{equation*}
\tfrac{1}{2}\Big[\tfrac{1}{2}n(n-1)\Big]\Big[\tfrac{1}{2}n(n-1)+1\Big] = \tfrac{1}{8}\big(n^4 - 2n^3 + 3n^2 - 2n\big) \tag{3.134}
\end{equation*}

个独立分量。我们还必须处理附加的对称性 (3.133)。(3.133) 的一个直接推论是，Riemann 张量的全反对称（totally antisymmetric）部分为零，

\begin{equation*}
R_{[\rho\sigma\mu\nu]} = 0. \tag{3.135}
\end{equation*}

事实上，这个方程再加上其他对称性 (3.129)、(3.130) 和 (3.131)，就足以蕴含 (3.133)；只需展开 (3.135) 并处理所得各项即可轻易证明。因此，一旦其他对称性都已计入，附加约束 (3.135) 就等价于 (3.133)。这代表多少个独立的限制？让我们设想分解

\begin{equation*}
R_{\rho\sigma\mu\nu} = X_{\rho\sigma\mu\nu} + R_{[\rho\sigma\mu\nu]}. \tag{3.136}
\end{equation*}

容易看出，任何全反对称的 4 指标张量都自动在其第一个和最后一个指标上反对称，并且在两对指标互换下对称。因此这些性质是对 $X_{\rho\sigma\mu\nu}$ 的独立限制，与要求 (3.135) 无关。一个全反对称的 4 指标张量有 $n(n-1)(n-2)(n-3)/4!$ 项，因此 (3.135) 把独立分量的数目削减了这么多。剩下

\begin{equation*}
\tfrac{1}{8}\big(n^4 - 2n^3 + 3n^2 - 2n\big) - \tfrac{1}{24}n(n-1)(n-2)(n-3) = \tfrac{1}{12}n^2\big(n^2 - 1\big) \tag{3.137}
\end{equation*}

个独立的 Riemann 张量分量。

因此在四维情形，Riemann 张量有 20 个独立分量。（在一维时它一个也没有。）这二十个函数恰好对应我们在第 2 章初次讨论局部惯性坐标时、无法靠巧妙选取坐标而消去的度规二阶导数的 20 个自由度。这应当会增强你的信心：Riemann 张量是曲率的恰当量度。

除了 Riemann 张量的代数对称性（它们约束任一点上独立分量的数目）之外，它还服从一个微分恒等式，约束它在不同点上的相对取值。考虑 Riemann 张量的协变导数，在局部惯性坐标中取值：

\begin{align*}
\nabla_{\hat\lambda}R_{\hat\rho\hat\sigma\hat\mu\hat\nu} &= \partial_{\hat\lambda}R_{\hat\rho\hat\sigma\hat\mu\hat\nu}\\
&= \tfrac{1}{2}\partial_{\hat\lambda}\big(\partial_{\hat\mu}\partial_{\hat\sigma}g_{\hat\rho\hat\nu} - \partial_{\hat\mu}\partial_{\hat\rho}g_{\hat\nu\hat\sigma} - \partial_{\hat\nu}\partial_{\hat\sigma}g_{\hat\rho\hat\mu} + \partial_{\hat\nu}\partial_{\hat\rho}g_{\hat\mu\hat\sigma}\big). \tag{3.138}
\end{align*}

对一个只在一点成立的表达式求导似乎不合规则，但我们忽略掉的项全都正比于 $\partial_{\hat\sigma}g_{\hat\mu\hat\nu}$，因而为零。我们想考虑头三个指标的循环置换之和：

\begin{align*}
&\nabla_{\hat\lambda}R_{\hat\rho\hat\sigma\hat\mu\hat\nu} + \nabla_{\hat\rho}R_{\hat\sigma\hat\lambda\hat\mu\hat\nu} + \nabla_{\hat\sigma}R_{\hat\lambda\hat\rho\hat\mu\hat\nu}\\
&= \tfrac{1}{2}\big(\partial_{\hat\lambda}\partial_{\hat\mu}\partial_{\hat\sigma}g_{\hat\rho\hat\nu} - \partial_{\hat\lambda}\partial_{\hat\mu}\partial_{\hat\rho}g_{\hat\nu\hat\sigma} - \partial_{\hat\lambda}\partial_{\hat\nu}\partial_{\hat\sigma}g_{\hat\rho\hat\mu} + \partial_{\hat\lambda}\partial_{\hat\nu}\partial_{\hat\rho}g_{\hat\mu\hat\sigma}\\
&\quad + \partial_{\hat\rho}\partial_{\hat\mu}\partial_{\hat\lambda}g_{\hat\sigma\hat\nu} - \partial_{\hat\rho}\partial_{\hat\mu}\partial_{\hat\sigma}g_{\hat\nu\hat\lambda} - \partial_{\hat\rho}\partial_{\hat\nu}\partial_{\hat\lambda}g_{\hat\sigma\hat\mu} + \partial_{\hat\rho}\partial_{\hat\nu}\partial_{\hat\sigma}g_{\hat\mu\hat\lambda}\\
&\quad + \partial_{\hat\sigma}\partial_{\hat\mu}\partial_{\hat\rho}g_{\hat\lambda\hat\nu} - \partial_{\hat\sigma}\partial_{\hat\mu}\partial_{\hat\lambda}g_{\hat\nu\hat\rho} - \partial_{\hat\sigma}\partial_{\hat\nu}\partial_{\hat\rho}g_{\hat\lambda\hat\mu} + \partial_{\hat\sigma}\partial_{\hat\nu}\partial_{\hat\lambda}g_{\hat\mu\hat\rho}\big)\\
&= 0. \tag{3.139}
\end{align*}

再说一次，既然这是张量之间的方程，尽管我们是在一个特定坐标系中导出它的，它在任何坐标系中都成立。此刻我们已经能够认出，反对称性 $R_{\rho\sigma\mu\nu} = -R_{\sigma\rho\mu\nu}$ 使我们可以把这个结果写成

\begin{equation*}
\nabla_{[\lambda}R_{\rho\sigma]\mu\nu} = 0. \tag{3.140}
\end{equation*}

这就是所谓的 \textbf{Bianchi 恒等式}（Bianchi identity）。对于一般的联络，还会有涉及挠率张量的附加项。它与 Jacobi 恒等式（Jacobi identity）密切相关，因为（回忆用协变导数的对易子给出的 Riemann 张量的定义）它表达的正是

\begin{equation*}
[[\nabla_\lambda, \nabla_\rho], \nabla_\sigma] + [[\nabla_\rho, \nabla_\sigma], \nabla_\lambda] + [[\nabla_\sigma, \nabla_\lambda], \nabla_\rho] = 0. \tag{3.141}
\end{equation*}

Riemann 张量有四个指标。有时，把一个张量表示成若干块之和是有用的，每一块各自更容易处理，并且可能有直接的物理解释。诀窍在于以坐标不变的方式做到这一点。例如，我们可以把 Riemann 张量分解为 $R^\rho{}_{\sigma ij}$ 和 $R^\rho{}_{\sigma i0}$，由它们可以重构整个张量（因为 $R^\rho{}_{\sigma 00}$ 为零）。但显然，这种分解在基的改变下不是不变的；我们要找的是在改变坐标时保持不变的分解。我们真正在做的事情，是考察洛伦兹群的表示。我们手头有两个基本技巧：取缩并，以及取对称/反对称部分。例如，给定任意一个 $(0,2)$ 型张量 $X_{\mu\nu}$，我们可以把它分解为对称和反对称部分，

\begin{equation*}
X_{\mu\nu} = X_{(\mu\nu)} + X_{[\mu\nu]}, \tag{3.142}
\end{equation*}

而对称部分还可以进一步分解为迹 $X = g^{\mu\nu}X_{(\mu\nu)}$ 和无迹部分 $\widehat X_{\mu\nu} = X_{(\mu\nu)} - \frac{1}{n}Xg_{\mu\nu}$，于是

\begin{equation*}
X_{\mu\nu} = \frac{1}{n}Xg_{\mu\nu} + \widehat X_{\mu\nu} + X_{[\mu\nu]}. \tag{3.143}
\end{equation*}

（注意 $X_{[\mu\nu]}$ 自动是无迹的。）当我们改变坐标时，各个部分 $Xg_{\mu\nu}$、$\widehat X_{\mu\nu}$ 和 $X_{[\mu\nu]}$ 只转到自身，不会转到彼此；我们说它们定义了 $(0,2)$ 型张量空间中的“不变子空间”（invariant subspaces）。对于更复杂的张量，等价的分解可能就没这么简单了。

对 Riemann 张量，我们的第一步是取一个缩并，构成 \textbf{Ricci 张量}（Ricci tensor，又称里奇张量）：

\begin{equation*}
R_{\mu\nu} = R^\lambda{}_{\mu\lambda\nu}. \tag{3.144}
\end{equation*}

对于由任意（不必是 Christoffel 的）联络构成的曲率张量，可以取的独立缩并有若干个。我们主要关心的是 Christoffel 联络，对它而言 (3.144) 是唯一的独立缩并；其余缩并要么为零，要么与这个缩并相关。与 Christoffel 联络相联系的 Ricci 张量自动是对称的，

\begin{equation*}
R_{\mu\nu} = R_{\nu\mu}, \tag{3.145}
\end{equation*}

这是 Riemann 张量对称性的一个推论。Ricci 张量的迹称为 \textbf{Ricci 标量曲率}（Ricci scalar，或称\textbf{标量曲率}，curvature scalar）：

\begin{equation*}
R = R^\mu{}_\mu = g^{\mu\nu}R_{\mu\nu}. \tag{3.146}
\end{equation*}

我们也可以构造无迹部分 $\widehat R_{\mu\nu} = R_{\mu\nu} - \frac{1}{n}Rg_{\mu\nu}$，但事实证明它并不特别有用；更常见的做法是用 $R_{\mu\nu}$ 和 $R$ 来表达各个量。

Ricci 张量与 Ricci 标量包含了 Riemann 张量迹的全部信息，留给我们的就是无迹部分。这些由 \textbf{Weyl 张量}（Weyl tensor）刻画，它基本上就是去掉全部缩并之后的 Riemann 张量。在 $n$ 维时它由下式给出

\begin{align*}
C_{\rho\sigma\mu\nu} &= R_{\rho\sigma\mu\nu} - \frac{2}{n-2}\big(g_{\rho[\mu}R_{\nu]\sigma} - g_{\sigma[\mu}R_{\nu]\rho}\big)\\
&\quad + \frac{2}{(n-1)(n-2)}g_{\rho[\mu}g_{\nu]\sigma}R \tag{3.147}
\end{align*}

这个繁杂的公式经过精心设计，使得 $C_{\rho\sigma\mu\nu}$ 的一切可能缩并都为零，同时保持 Riemann 张量的对称性：

\begin{align*}
C_{\rho\sigma\mu\nu} &= C_{[\rho\sigma][\mu\nu]},\\
C_{\rho\sigma\mu\nu} &= C_{\mu\nu\rho\sigma},\\
C_{\rho[\sigma\mu\nu]} &= 0. \tag{3.148}
\end{align*}

Weyl 张量只在三维或更高维有定义，而在三维中它恒等于零。Weyl 张量最重要的性质之一，是它在共形变换（conformal transformation，见附录 G）下不变。这意味着，如果你对某个度规 $g_{\mu\nu}$ 计算出 $C^\rho{}_{\sigma\mu\nu}$（注意第一个指标在上），然后对由 $\omega^2(x)g_{\mu\nu}$ 给出的度规（其中 $\omega(x)$ 是时空上任意一个处处非零的函数）再计算一遍，你会得到相同的答案。因此它常被称为\emph{共形张量}（conformal tensor）。

Bianchi 恒等式的一个特别有用的形式来自对 (3.139) 缩并两次：

\begin{align*}
0 &= g^{\nu\sigma}g^{\mu\lambda}\big(\nabla_\lambda R_{\rho\sigma\mu\nu} + \nabla_\rho R_{\sigma\lambda\mu\nu} + \nabla_\sigma R_{\lambda\rho\mu\nu}\big)\\
&= \nabla^\mu R_{\rho\mu} - \nabla_\rho R + \nabla^\nu R_{\rho\nu}, \tag{3.149}
\end{align*}

或者

\begin{equation*}
\nabla^\mu R_{\rho\mu} = \tfrac{1}{2}\nabla_\rho R. \tag{3.150}
\end{equation*}

注意，与偏导数不同，由于度规相容性，对协变导数升指标是有意义的。我们定义 \textbf{Einstein 张量}（Einstein tensor）

% 未完，句子在书页 130 末中断，由下一块（书页 131 起）承接。
